% === Begin Label Data ===
% ID: 796
% 难度星级: 
% 标签: 
% 备注: 
% 组卷引用次数: 0
% === End  Label Data ===

\begin{problem}{2019}{G}{全国II卷}{5}{统计}
演讲比赛共有 $9$ 位评委分别给出某选手的原始评分，评定该选手的成绩时，从 $9$ 个原始评分中去掉 $1$ 个最高分、$1$ 个最低分，得到 $7$ 个有效评分．$7$ 个有效评分与 $9$ 个原始评分相比，不变的数字特征是 (\hspace{1cm})
\begin{choices}
\choice{{中位数}}
\choice{{平均数}}
\choice{{方差}}
\choice{{极差}}
\end{choices}
\end{problem}

\begin{answer}
A
\end{answer}

\begin{solutions}
将 $9$ 位评委分别给出的某选手的原始评分有小到大排列为 $a_1\leqslant a_2\leqslant\cdots\leqslant a_9$，则 $9$ 个原始评分的中位数为 $a_5$，平均数为 $\bar{a}=\dfrac{a_1+a_2+\cdots+a_9}{9}$，方差为 $\dfrac{1}{9}\displaystyle\sum_{i=1}^{9}(a_i-\bar{a})^2$，极差为 $a_9-a_1$．从 $9$ 个原始评分中去掉 $1$ 个最高分 $a_9$ 和 $1$ 个最低分 $a_1$，得到 $7$ 个有效评分，由上面四个数字特征的定义可知，只有中位数不受去掉 $1$ 个最高分和 $1$ 个最低分的影响，所以应该选择 A．
\end{solutions}